\documentclass[11pt]{article}
\usepackage[margin=1in]{geometry}
\usepackage{amsmath,amssymb,booktabs,siunitx,hyperref,microtype,xcolor}
\usepackage{enumitem}
\setlist[itemize]{leftmargin=*,topsep=2pt,itemsep=2pt}

\title{PNP--BV Forward Solver Evolution\\
       From 4-Species PNP to 3-Species + Boltzmann + Log-c + Log-Rate BV}
\author{Jake Weinstein}
\date{Week of April~27, 2026}

\newcommand{\true}{\mathrm{TRUE}}
\newcommand{\Eeq}{E_{\mathrm{eq}}}
\newcommand{\VRHE}{V_{\mathrm{RHE}}}
\newcommand{\VT}{V_{T}}

\begin{document}
\maketitle

%% ============================================================
\section{Summary}
%% ============================================================
The forward solver supporting the 4-parameter inverse problem
($\log k_{0,1}, \log k_{0,2}, \alpha_1, \alpha_2$) was rebuilt over the past
two weeks. Three independent changes, each motivated by a distinct numerical
failure of the original 4-species PNP solver, are now stacked in the active
configuration:
\begin{enumerate}
  \item \textbf{4-species $\to$ 3-species PNP + analytic Boltzmann counterion}
        --- removes the inert supporting-electrolyte counterion as a dynamic
        Nernst--Planck unknown while retaining its contribution to Poisson's
        equation.
  \item \textbf{Concentration $c_i \to$ log-concentration $u_i = \ln c_i$}
        --- enforces positivity of the dynamic species and avoids the negative
        concentration iterates that previously triggered floor/clip
        regularizations.
  \item \textbf{Standard Butler--Volmer $\to$ log-rate Butler--Volmer}
        --- removes a phantom $\mathrm{R}_2$ sink at high anodic overpotential
        caused by clipping the surface concentration before evaluating the
        boundary reaction rate.
\end{enumerate}
The combined effect is that the usable voltage grid expands from roughly
$\VRHE \le \SI{+0.20}{\volt}$ to $\VRHE \le \SI{+0.60}{\volt}$. In the local
Fisher analysis at TRUE parameter values, the old nearly singular
$\log k_{0,2}$ direction is removed: $\mathrm{cond}(F)$ improves from about
$2\times 10^{11}$ to about $1.79\times 10^{7}$, and the canonical-ridge cosine
drops from $1.000$ to $0.031$. This is a major improvement, but it should not
be read as a complete solution of the inverse problem. Clean-data inverse runs
still show multiple basins: $\alpha_1$ and $\alpha_2$ are recovered to
$\le 2\%$ in 2/4 standard initializations and to $\le 16\%$ in 3/4, while no
single initialization recovers all four parameters jointly.

%% ============================================================
\section{Original 4-Species Solver}
%% ============================================================
The original forward model was a coupled Poisson--Nernst--Planck (PNP)
system on a graded 2D rectangular mesh with four species:
$\mathrm{O}_2$, $\mathrm{H}_2\mathrm{O}_2$, $\mathrm{H}^{+}$, and the
supporting-electrolyte counterion $\mathrm{ClO}_4^{-}$. In dimensional form,
the bulk equations for each species $i$ are
\begin{equation}
  \frac{\partial c_i}{\partial t}
  \;=\;
  -\,\nabla\!\cdot J_i,
  \qquad
  J_i = -D_i\left(\nabla c_i + \frac{z_i F}{RT}c_i\nabla\varphi\right),
  \qquad
  -\,\varepsilon\,\nabla^{2}\varphi \;=\; F\sum_i z_i c_i,
  \label{eq:pnp-orig}
\end{equation}
with Butler--Volmer (BV) flux boundary conditions at the electrode for the
two reactions
\begin{align}
  \mathrm{R}_1\!&: \;\mathrm{O}_2 + 2\mathrm{H}^{+} + 2e^{-}
                       \to \mathrm{H}_2\mathrm{O}_2,
  \qquad &\Eeq^{(1)} = \SI{0.68}{\volt}, \\
  \mathrm{R}_2\!&: \;\mathrm{H}_2\mathrm{O}_2 + 2\mathrm{H}^{+} + 2e^{-}
                       \to 2\mathrm{H}_2\mathrm{O},
  \qquad &\Eeq^{(2)} = \SI{1.78}{\volt}.
\end{align}
The implementation uses the corresponding nondimensional form, where voltage
is scaled by $\VT = RT/F$.

\paragraph{Where it failed.}
\begin{itemize}
  \item \textbf{Voltage ceiling near $\VRHE \approx \SI{+0.20}{\volt}$.} At
        higher anodic potentials the $\mathrm{H}_2\mathrm{O}_2$ surface
        concentration can drop by many orders of magnitude. Newton iterates in
        the concentration formulation could become non-positive, triggering
        floor regularizations whose Jacobian contributions stalled the SNES
        solve.
  \item \textbf{Nearly singular FIM at TRUE.} Restricted to
        $\VRHE \le 0.20$\,V, the CD+PC Fisher information had condition number
        $\approx 2\times 10^{11}$ with weak eigenvector almost pure
        $\log k_{0,2}$. In that voltage window, $k_{0,2}$ was practically
        unidentifiable.
  \item \textbf{$\mathrm{ClO}_4^{-}$ added cost without chemistry.} The
        counterion is inert in this model and exists to provide charge balance.
        Solving for it as a fourth dynamic NP variable increased assembly and
        nonlinear-solve cost while not adding an electrode reaction or an
        independent inverse observable.
\end{itemize}

%% ============================================================
\section{Change~1: 3-Species PNP + Analytic Boltzmann Counterion}
%% ============================================================
\paragraph{Motivation.}
$\mathrm{ClO}_4^{-}$ is inert and present as supporting electrolyte. Rather
than solving a fourth dynamic Nernst--Planck equation for this nonreactive
counterion, the reduced model assumes it remains in local electrochemical
equilibrium with the electrostatic potential. Under that assumption its
profile is represented analytically by a Boltzmann factor,
\(
  c_{\mathrm{ClO}_4}(x) = c_{\mathrm{bulk}}\exp\!\bigl(-z_{\mathrm{ClO}_4}\varphi(x)/\VT\bigr)
\),
with the sign following the code's potential convention. This is a modeling
reduction for an inert supporting electrolyte, not a new physical mechanism.

\paragraph{Implementation.}
Drop $\mathrm{ClO}_4^{-}$ as a primary unknown. The Poisson equation keeps its
analytic contribution to the charge density:
\begin{equation}
  -\,\varepsilon\nabla^{2}\varphi
   \;=\; F\!\!\sum_{i\in\{\mathrm{O}_2,\,\mathrm{H}_2\mathrm{O}_2,\,\mathrm{H}^{+}\}} z_i c_i
   \;-
   F\,c_{\mathrm{ClO}_4,\mathrm{bulk}}\,\exp(-\varphi/\VT),
\end{equation}
where the exponential sign should be interpreted consistently with the
nondimensional potential convention used in the code. The three remaining
species have charges $z = [0,0,+1]$ and retain the full NP coupling.
Operationally this is implemented through the Boltzmann contribution to the
Poisson residual; see \texttt{add\_boltzmann()} in the current study scripts.

\paragraph{Result.}
The reduced model removes one dynamic species and produced a measurable speedup
in the forward studies. To quantify how closely the reduced solver tracks
the full 4-species PNP, the current production 3-species $+$ Boltzmann log-c
forward solver (the build path used by \texttt{v18\_logc\_lsq\_inverse.py},
the most recent inverse) was run head-to-head against the 4-species standard
solver via \texttt{solve\_grid\_with\_charge\_continuation} on a shared mesh,
shared TRUE kinetic parameters, and shared $E_{\mathrm{eq}}$. The 4-species
charge-continuation pipeline reaches $z=1$ across the documented window
$\VRHE \in [-0.50, +0.10]$\,V, but the 3-species log-c solver only
cold-starts in $\VRHE \gtrsim -0.20$\,V. To grow the overlap set we add a
voltage-continuation strategy on the 3-species side: cold-start where
possible, then warm-start at full charge from the nearest converged
neighbour and march the applied potential outward in $\Delta V$ substeps
with bisection fallback. With four substeps per
$\SI{0.10}{\volt}$ jump, this fills the entire $\VRHE \in [-0.50, +0.10]$\,V
grid. Across the resulting eight-voltage overlap set, both observables
agree well within the F2 diffusion-limit tolerance ($5\%$ of
$\max|\text{obs}|$) that the inverse-study validation framework uses
elsewhere (\texttt{Forward/bv\_solver/validation.py}):
\begin{center}
  \begin{tabular}{lcc}
    \toprule
                                       & \textbf{max} & \textbf{mean} \\
    \midrule
    $|\Delta \mathrm{CD}| / \max|\mathrm{CD}|$ & $0.269\%$ & $0.201\%$ \\
    $|\Delta \mathrm{PC}| / \max|\mathrm{PC}|$ & $0.104\%$ & $0.033\%$ \\
    \bottomrule
  \end{tabular}
\end{center}
All eight overlap points pass the $5\%$ F2 threshold on both observables.
At anodic voltages where peroxide current is small ($\VRHE \ge -0.10$\,V),
PC differences appear as a sign flip on values four to seven orders of
magnitude below $\max|\mathrm{PC}|\approx 1.8\times 10^{-1}$ --- a numerical
artefact of evaluating a near-zero quantity in two different formulations,
not a physics disagreement. The important claim is therefore not that the
two formulations are algebraically identical, but that this controlled
reduction preserves the ORR observables to better than half a percent of
the diffusion-limit scale across the entire window where the 4-species
solver itself converges --- well inside the accuracy needed for the current
inverse-study comparisons. Raw data, CSV, plot, and summary are saved under
\path{StudyResults/v24_3sp_logc_vs_4sp_validation/}; the comparison script
(including the warm-start continuation strategy) is
\path{scripts/studies/v24_3sp_logc_vs_4sp_validation.py}.

%% ============================================================
\section{Change~2: Log-Concentration Formulation}
%% ============================================================
\paragraph{Motivation.}
Concentrations near a reactive electrode can span many orders of magnitude. In
the diffuse layer, $c_i(x)$ varies approximately as
$c_{\mathrm{bulk}}\exp(-z_i\varphi/\VT)$. When $c$ is the primary unknown, the
SNES line search can overshoot near small concentrations and produce
non-positive iterates. Those iterates then require floors or clips, and the
regularization can contaminate the Jacobian and stall continuation.

\paragraph{Implementation.}
Switch to $u_i = \ln c_i$ as the primary unknown. The Nernst--Planck weak form
becomes
\begin{equation}
  \int_\Omega
    \frac{e^{u_i} - e^{u_i^{\mathrm{old}}}}{\Delta t}\,v\;\mathrm{d}x
  \;+\;
  \int_\Omega
    D_i\,e^{u_i}\bigl(\nabla u_i + z_i \nabla\phi\bigr)\!\cdot\!\nabla v\;\mathrm{d}x
  \;=\;
  \text{(BV boundary terms)},
\end{equation}
and Poisson's right-hand side replaces $c_i$ by $\exp(u_i)$. See
\texttt{Forward/bv\_solver/forms\_logc.py}.

\paragraph{Why it helps.}
\begin{itemize}
  \item $c_i = \exp(u_i) > 0$ by construction, so the dynamic species remain
        positive in the bulk formulation.
  \item The Debye-layer behavior becomes approximately linear in $u_i$:
        $u_i(x) \approx \ln c_{\mathrm{bulk}} - z_i\varphi/\VT$, which is more
        naturally represented on the existing graded mesh.
  \item Remaining numerical guards/clamps should be understood as overflow
        protection rather than concentration floors in the physical variable;
        later clamp sweeps checked that these guards did not drive the reported
        Fisher results.
\end{itemize}

%% ============================================================
\section{Change~3: Log-Rate Butler--Volmer}
%% ============================================================
\paragraph{Starting point.}
The reference for our group's canonical BV expression is the standard
two-branch form for an $n$-electron redox couple
$\mathrm{O} + n e^{-} \rightleftharpoons \mathrm{R}$%
\footnote{See the internal document \texttt{docs/PNP Equation
Formulations.tex} for the group's full derivation, weak forms, and
flux-boundary mapping. The form below is its
$R_{\mathrm{BV}}$ expression.}:
\begin{equation}
  \label{eq:bv-canonical}
  R_{\mathrm{BV}}
  \;=\;
  k_0
  \!\left[\,
  c_{\mathrm{O}}\,\exp\!\Bigl(-\alpha\,\tfrac{n F \eta}{RT}\Bigr)
  \;-\;
  c_{\mathrm{R}}\,\exp\!\Bigl((1-\alpha)\,\tfrac{n F \eta}{RT}\Bigr)
  \,\right],
\end{equation}
where $c_{\mathrm{O}}, c_{\mathrm{R}}$ are the surface concentrations on the
electrode reaction plane and $\eta = \phi_m - \phi_s - \Eeq$ is the
overpotential. In the nondimensional code, $RT/F$ is replaced by $\VT$ and
$n F \eta / RT$ becomes $n\eta/\VT$. Equation~\eqref{eq:bv-canonical} is
the algebraic object both formulations evaluate; the two formulations
disagree only on \emph{how} the surface contribution is combined with the
voltage exponential.

\paragraph{Two ways to evaluate it in a log-c solver.}
With $u_i = \ln c_i$ as the primary unknown, $c_{\mathrm{O}} = \exp(u_{\mathrm{O}})$
exactly. The two evaluations differ only at the boundary, in how $u_i$ is
turned back into a rate. Showing the cathodic branch only for clarity:
\begin{center}
\renewcommand{\arraystretch}{1.6}
\begin{tabular}{@{}p{0.46\textwidth}@{\hspace{0.04\textwidth}}p{0.46\textwidth}@{}}
\toprule
\textbf{Old: clip-then-multiply} & \textbf{New: log-rate} \\
\midrule
\(c_{\mathrm{surf},i} = \exp\!\bigl(\mathrm{clip}(u_i,\pm 30)\bigr)\)
&
(no surface clip)
\\
\(r_j \;=\; k_{0,j}\;
            c_{\mathrm{surf},\mathrm{cat}(j)}\;
            \exp(-\alpha_j\,n_e\,\eta_j/\VT)\)
&
\(\log r_j \;=\;
   \ln k_{0,j} \;+\; u_{\mathrm{cat}(j)}
                  \;-\; \alpha_j\,n_e\,\eta_j/\VT\)
\\
&
\(r_j \;=\; \exp(\log r_j)\)
\\
\midrule
exponentiate $u$, clip, multiply
&
add $u$ inside the exponent, exponentiate once
\\
\bottomrule
\end{tabular}
\end{center}
The anodic branch is analogous: replace $u_{\mathrm{cat}(j)}$ by
$u_{\mathrm{anod}(j)}$ and the sign of the $\eta$ term, then subtract the
two exponentials. The full implementation also adds the
$\sum_f \nu_f (u_{\mathrm{sp}(f)} - \ln c_{\mathrm{ref},f})$ contribution
from cathodic concentration factors (e.g.\ the H$^+$ stoichiometric power
in our R$_1$/R$_2$ reactions); under the new form those terms also enter
additively inside the exponent.

\paragraph{Why the old form fails at high $\eta$.}
The two formulations agree algebraically when $|u_i|$ stays modest. They
disagree when $u_i$ runs to large negative values --- exactly the regime
this work was trying to access. In the old form, $u_{\mathrm{H}_2\mathrm{O}_2}$
should tend to a very negative number at high anodic $\VRHE$, but the
$\pm 30$ clip pins $c_{\mathrm{surf,H_2O_2}}$ at a finite floor. Multiplying
that floor by the large $\mathrm{R}_2$ exponential produces a phantom
cathodic sink. In the new form, the same large-negative $u_i$ is added
\emph{inside} the exponent, so $\log r_j \to -\infty$ and $r_j \to 0$
smoothly, which is the physically correct behavior. This was a major source
of Newton failures in the $\VRHE \in [+0.30,+0.60]$\,V range; the
log-rate form removes them. The toggle is
\texttt{bv\_log\_rate=True} in the BV convergence configuration; see
\texttt{forms\_logc.py:294--360}.

\paragraph{Result.}
The extended voltage grid
$\VRHE \in \{-0.10, +0.10, +0.20, +0.30, +0.40, +0.50, +0.60\}$\,V now
converges under the current cold-ramp/charge-continuation pipeline. Some
points require the continuation machinery rather than a single direct Newton
solve, so the safe statement is ``converges under the current pipeline,'' not
``converges trivially from every cold start.''

%% ============================================================
\section{Combined Results}
%% ============================================================

\subsection*{Voltage window}
Pre-rebuild, the usable synthetic voltage window was roughly
$\VRHE \le \SI{+0.20}{\volt}$. Post-rebuild, the current pipeline reaches the
extended grid through $\VRHE = \SI{+0.60}{\volt}$. This matters because
$\mathrm{R}_2$ begins to unclip at approximately
\(
  \VRHE > \Eeq^{(2)} - 50\VT \approx \SI{+0.495}{\volt},
\)
under the code's $\eta_{\mathrm{scaled}}$ clipping convention. The inverse
problem now has access to voltage-shape information from a regime that was
previously unreachable.

\subsection*{Fisher information}
At TRUE parameter values, with current density (CD) and peroxide current (PC)
as observables and a global-max $2\%$ noise model:
\begin{center}
  \begin{tabular}{lcc}
    \toprule
                                       & \textbf{old grid}    & \textbf{new grid} \\
                                       & ($\le \SI{+0.20}{V}$) & ($\le \SI{+0.60}{V}$) \\
    \midrule
    $\sigma_{\min}(S_{\mathrm{white}})$ & $2.35\times 10^{-2}$ & $2.51$ \\
    $\mathrm{cond}(F)$                 & $2.0\times 10^{11}$  & $1.79\times 10^{7}$ \\
    canonical-ridge $|\cos|$           & $1.000$              & $0.031$ \\
    weak eigenvector                   & nearly pure $\log k_{0,2}$  & almost pure $\log k_{0,1}$ \\
    \bottomrule
  \end{tabular}
\end{center}
The historic ``$\log k_{0,2}$ unidentifiable'' local Fisher result is gone.
A new weak direction in $\log k_{0,1}$ remains. Thus the local FIM is much
better conditioned than before, but it is still not perfectly conditioned, and
it does not by itself guarantee global recovery from arbitrary initializations.

\subsection*{Adjoint validity}
At three of the previously inaccessible voltages ($\VRHE \in
\{0.30,\,0.40,\,0.50\}$), an initial adjoint-vs-FD check appeared to fail by
exactly a factor of two. Detailed sweeps over annotated SNES iteration count,
steady-state tolerance, FD step size, and cold-ramp versus warm-start FD showed
that the discrepancy came from the warm-start FD calculation, not from the
adjoint. With cold-ramp finite differences, the relevant nonzero sensitivity
components agree with the adjoint to about six significant figures.

\subsection*{Inverse-problem implications}
With the new forward solver, the four-parameter trust-region inverse with
clean data and four standard initializations produces:
\begin{itemize}
  \item $\alpha_1, \alpha_2$ recovered to $\le 2\%$ in 2/4 initializations and
        to $\le 16\%$ in 3/4 initializations;
  \item each individual $k_{0,j}$ recovered to $\le 5\%$ from at least one
        initialization;
  \item but no single initialization recovers all four parameters jointly.
\end{itemize}
The $(k_0,\alpha)$ Tafel-ridge structure is no longer a single-direction local
singularity dominated by $\log k_{0,2}$. The remaining problem is a multi-basin
global cost surface: the data constrain some kinetic combinations strongly,
but the optimizer can still land in different Tafel-like basins depending on
initialization.

%% ============================================================
\section{What This Does and Does Not Solve}
%% ============================================================
\paragraph{Solved or substantially improved.}
\begin{itemize}
  \item Forward simulation now reaches the extended synthetic voltage grid
        through $\SI{+0.60}{\volt}$ under the current continuation pipeline.
  \item The old local Fisher weak direction dominated by $\log k_{0,2}$ is
        removed once log-rate BV and the extended grid are used.
  \item $\alpha_1$ and $\alpha_2$ are substantially more robustly determined
        than in the earlier single-experiment formulation, although recovery is
        still initialization-dependent.
\end{itemize}

\paragraph{Not solved.}
\begin{itemize}
  \item $\mathrm{cond}(F) \approx 1.79\times 10^7$ is a major improvement over
        the previous $\approx 2\times 10^{11}$, but it is not numerically
        ``perfectly conditioned.''
  \item The global cost surface for the four-parameter joint inverse contains
        multiple Tafel-ridge basins; the optimizer's landing basin depends on
        initialization.
  \item Soft priors at literature-defensible strength
        ($\sigma_{\log k_0} = \ln 3$ or $\ln 10$) shift within-basin position
        but do not bridge the inter-basin barriers observed in the current
        clean-data studies.
\end{itemize}
For the current inverse project, the forward solver is no longer the obvious
primary bottleneck. The next-step plan is therefore to close out the
single-experiment route with restart diagnostics and anchored Tafel-coordinate
optimization, then move to multi-experiment Fisher-information design, starting
with bulk-$\mathrm{O}_2$ variation.

%% ============================================================
\section{Relation to Prior Work}
%% ============================================================
None of the three changes is methodologically novel in isolation; the
contribution here is the engineering combination and its quantitative effect
on Fisher conditioning for the ORR inverse problem. The relevant precedents,
to the extent the literature was surveyed:

\paragraph{Log-concentration / log-density formulations (Change~2).}
The substitution $u_i = \ln c_i$ has a long history in adjacent fields. In
semiconductor drift--diffusion, the closely related \emph{Slotboom} (or
quasi-Fermi-level) variables have been standard since the late 1960s, and
the Scharfetter--Gummel finite-volume discretization built around them
remains a default in device simulation~\cite{ScharfetterGummel1969}. For the
Poisson--Nernst--Planck system specifically, Metti, Xu, and Liu introduced
what they explicitly named the ``log-density formulation'' as a means of
constructing positivity-preserving and energetically stable
schemes~\cite{MettiXuLiu2016}. Recent extensions push this idea to high-order
space--time finite element methods with arbitrary order positivity and
unconditional energy stability~\cite{LiuMaxianShenZitelliWaltzMonteiro2022}.
Our usage applies this established idea inside a Firedrake-based ORR
PNP--BV solver; we do not claim novelty in the substitution itself.

\paragraph{Reduced supporting-electrolyte models (Change~1).}
Replacing the inert supporting electrolyte by an analytic Boltzmann profile
in the Poisson equation while keeping the reactive species as full
Nernst--Planck unknowns is a known model reduction commonly called the
\emph{Poisson--Boltzmann--Nernst--Planck} (PBNP) hybrid
model~\cite{ZhengChenWei2011}. It is used extensively for ion-channel and
biomolecular simulations, where most ionic species are passive bath and only
a few are chemically active. Our 3sp-plus-Boltzmann setup adapts that
reduction to a two-step ORR system; the modeling assumption (inert,
locally-equilibrated supporting electrolyte) is the same.

\paragraph{Log-form rate evaluation (Change~3).}
Carrying chemical or electrochemical rates in log space to avoid floating
overflow or underflow is standard in stiff combustion and plasma chemistry
codes, where Arrhenius rates routinely span tens of orders of magnitude. In
the electrochemistry literature, the well-known logarithmic form of the
Butler--Volmer relation is the \emph{Tafel} approximation~\cite{Tafel1905},
which drops one of the two BV branches in the high-overpotential limit and
is therefore mathematically distinct from the substitution used here. The
present writeup retains both anodic and cathodic branches, but evaluates
each branch as $r_j = \exp(\log r_j)$ inside the boundary residual so that
the unclamped log-concentrations enter additively. We did not locate a
specific cited precedent for that exact construction in the BV
electrochemistry literature; it may be a small local numerical novelty,
although the underlying log-form-rate idea is not.

\paragraph{Tafel approximation versus log-rate BV.}
To make the distinction precise: the Tafel form is a \emph{physical
approximation} valid only at large $|\eta|$, obtained by dropping the
exponentially smaller BV branch; the log-rate BV evaluation is an
\emph{algebraic identity} that keeps both branches at every $\eta$,
including the cross-over regime $\eta \to 0$ where Tafel is invalid. Both
forms incidentally avoid the boundary-level surface-concentration clip at
high $|\eta|$---Tafel because only one exponential remains, log-rate BV
because the relevant $u_i$ enters additively before exponentiation---but
the two are not interchangeable. A forward PNP--BV solver that must run
across a voltage grid containing $\eta \approx 0$ transitions needs the
identity, not the asymptote. Conversely, when fitting a Tafel slope from
experimental log-current data, the Tafel approximation is the appropriate
construction; log-rate BV would be the wrong layer of the modeling
stack.

%% ============================================================
\section*{Code Pointers}
%% ============================================================
\begingroup
\small
\sloppy
\begin{itemize}
  \item Log-c PNP weak forms: \path{Forward/bv_solver/forms_logc.py}. Key
        functions include \texttt{build\_context\_logc},
        \texttt{build\_forms\_logc}, and
        \texttt{set\_initial\_conditions\_logc}.
  \item Log-rate BV toggle: \texttt{bv\_log\_rate} in
        \path{Forward/bv_solver/config.py}; consumed in
        \path{forms_logc.py}, around lines 294--360.
  \item Boltzmann counterion residual: \texttt{add\_boltzmann()} in the
        top-level study scripts, including
        \path{scripts/studies/v18_logc_lsq_inverse.py}.
  \item Charge-continuation pipeline:
        \path{Forward/bv_solver/grid_charge_continuation.py}.
\end{itemize}
\endgroup

%% ============================================================
\appendix
\section{Per-voltage 4sp vs 3sp+Boltzmann log-c overlap data}
%% ============================================================
Per-voltage observable values from
\path{StudyResults/v24_3sp_logc_vs_4sp_validation/comparison.csv}, generated
by \path{scripts/studies/v24_3sp_logc_vs_4sp_validation.py}. Both solvers
share the same graded mesh (Nx$=8$, Ny$=200$, $\beta=3$), the same TRUE
kinetic parameters ($\hat{k}_{0,1} = 1.263158\times 10^{-3}$,
$\hat{k}_{0,2} = 5.263158\times 10^{-5}$, $\alpha_1 = 0.627$,
$\alpha_2 = 0.5$), the same equilibrium potentials ($E_{\mathrm{eq}}^{(1)} =
\SI{0.68}{\volt}$, $E_{\mathrm{eq}}^{(2)} = \SI{1.78}{\volt}$ vs RHE), and
the same $I_{\mathrm{SCALE}}$. The 4-species side runs through
\texttt{solve\_grid\_with\_charge\_continuation}; the 3-species log-c side
runs the production build path from \texttt{v18\_logc\_lsq\_inverse.py},
cold-starting where possible and warm-starting (with paf substepping plus
bisection) from the nearest converged neighbour everywhere else. Errors
are reported relative to $\max|\mathrm{obs}|$ on the eight-voltage grid:
$\max|\mathrm{CD}| \approx 1.84\times 10^{-1}$,
$\max|\mathrm{PC}| \approx 1.81\times 10^{-1}$ (in nondimensional flux units;
both solvers convert to mA/cm$^{2}$ via $-I_{\mathrm{SCALE}}$).

\begingroup
\small
\setlength{\tabcolsep}{4pt}
\renewcommand{\arraystretch}{1.05}
\begin{center}
\begin{tabular}{r r r r r r r l c}
\toprule
$\VRHE$ (V) & $\mathrm{CD}_{\mathrm{4sp}}$ & $\mathrm{CD}_{\mathrm{3sp}}$
            & $|\Delta \mathrm{CD}|/\mathrm{CD}_{\max}\,\%$
            & $\mathrm{PC}_{\mathrm{4sp}}$ & $\mathrm{PC}_{\mathrm{3sp}}$
            & $|\Delta \mathrm{PC}|/\mathrm{PC}_{\max}\,\%$
            & 3sp method & verdict \\
\midrule
$-0.500$ & $-1.8390\mathrm{e}{-1}$ & $-1.8410\mathrm{e}{-1}$ & $0.108$
         & $-1.8112\mathrm{e}{-1}$ & $-1.8131\mathrm{e}{-1}$ & $0.104$
         & warm$\leftarrow{-0.40}$ & PASS \\
$-0.400$ & $-1.8283\mathrm{e}{-1}$ & $-1.8305\mathrm{e}{-1}$ & $0.122$
         & $-1.6374\mathrm{e}{-1}$ & $-1.6391\mathrm{e}{-1}$ & $0.094$
         & warm$\leftarrow{-0.30}$ & PASS \\
$-0.300$ & $-1.8133\mathrm{e}{-1}$ & $-1.8172\mathrm{e}{-1}$ & $0.212$
         & $-6.4999\mathrm{e}{-2}$ & $-6.5039\mathrm{e}{-2}$ & $0.022$
         & warm$\leftarrow{-0.20}$ & PASS \\
$-0.200$ & $-1.7968\mathrm{e}{-1}$ & $-1.8017\mathrm{e}{-1}$ & $0.268$
         & $-1.5838\mathrm{e}{-3}$ & $-1.5678\mathrm{e}{-3}$ & $0.009$
         & cold & PASS \\
$-0.100$ & $-1.7753\mathrm{e}{-1}$ & $-1.7802\mathrm{e}{-1}$ & $0.266$
         & $-1.2483\mathrm{e}{-5}$ & $+2.9058\mathrm{e}{-6}$ & $0.008$
         & cold & PASS \\
$+0.000$ & $-1.7330\mathrm{e}{-1}$ & $-1.7379\mathrm{e}{-1}$ & $0.266$
         & $-9.6937\mathrm{e}{-8}$ & $+1.5288\mathrm{e}{-5}$ & $0.008$
         & cold & PASS \\
$+0.050$ & $-1.6936\mathrm{e}{-1}$ & $-1.6985\mathrm{e}{-1}$ & $0.269$
         & $-8.6209\mathrm{e}{-9}$ & $+1.5378\mathrm{e}{-5}$ & $0.008$
         & cold & PASS \\
$+0.100$ & $-1.6289\mathrm{e}{-1}$ & $-1.6307\mathrm{e}{-1}$ & $0.100$
         & $-7.7618\mathrm{e}{-10}$ & $+1.5388\mathrm{e}{-5}$ & $0.008$
         & cold & PASS \\
\midrule
\multicolumn{3}{r}{\textbf{aggregate}} & $\mathbf{0.269}$ (max)
         & & & $\mathbf{0.104}$ (max) & & 8/8 PASS \\
\multicolumn{3}{r}{} & $0.201$ (mean) & & & $0.033$ (mean) & & \\
\bottomrule
\end{tabular}
\end{center}
\endgroup

\noindent\textbf{Reading the table.} CD agreement stays under $0.27\%$ of
$\mathrm{CD}_{\max}$ across the full eight-voltage grid, including the
three cathodic points reached only by warm-start continuation. PC
agreement stays under $0.11\%$ of $\mathrm{PC}_{\max}$ everywhere. The
sign flips on PC at $\VRHE \ge -0.10$\,V correspond to PC values four to
seven orders of magnitude below $\mathrm{PC}_{\max}$ --- a numerical
artefact of evaluating a near-zero quantity in two different formulations,
not a physics disagreement. Wall time on this run was 17.8\,s for the
4-species sweep and 74.3\,s for the 3-species cold/warm pipeline.

%% ============================================================
\begin{thebibliography}{9}
%% ============================================================

\bibitem{ScharfetterGummel1969}
Scharfetter, D.\,L., and Gummel, H.\,K.,
\emph{Large-signal analysis of a silicon Read diode oscillator},
IEEE Trans.\ Electron Devices, \textbf{16}(1), 64--77 (1969).
The original drift-diffusion finite-volume discretization built around
Slotboom-type variables; widely cited as the canonical positivity-preserving
scheme for charge-transport problems mathematically equivalent to PNP.

\bibitem{Tafel1905}
Tafel, J.,
\emph{\"Uber die Polarisation bei kathodischer Wasserstoffentwicklung},
Z.\ Phys.\ Chem.\ \textbf{50}, 641--712 (1905).
Original derivation of the high-overpotential logarithmic relation between
overpotential and current that bears Tafel's name; mathematically distinct
from the log-rate BV evaluation used here, but the historical antecedent for
``log-form'' relations in electrochemistry.

\bibitem{ZhengChenWei2011}
Zheng, Q., Chen, D., and Wei, G.\,W.,
\emph{Second-order Poisson--Nernst--Planck solver for ion transport},
J.\ Comput.\ Phys.\ \textbf{230}(13), 5239--5262 (2011);
and the closely related ``Poisson--Boltzmann--Nernst--Planck model,''
J.\ Chem.\ Phys.\ \textbf{134}, 194101 (2011),
\href{https://doi.org/10.1063/1.3581031}{doi:10.1063/1.3581031}.
Establishes the PBNP hybrid: dynamic NP for chemically active species,
Boltzmann factor for passive supporting electrolyte.

\bibitem{MettiXuLiu2016}
Metti, M.\,S., Xu, J., and Liu, C.,
\emph{Energetically stable discretizations for charge transport and
electrokinetic models},
J.\ Comput.\ Phys.\ \textbf{306}, 1--18 (2016),
\href{https://doi.org/10.1016/j.jcp.2015.10.053}{doi:10.1016/j.jcp.2015.10.053}.
Introduces the substitution $u_i = \ln c_i$ for PNP under the explicit name
\emph{log-density formulation}, with positivity preservation and
free-energy dissipation as the headline numerical properties.

\bibitem{LiuMaxianShenZitelliWaltzMonteiro2022}
Liu, H., Maxian, O., Shen, J., Zitelli, J., Waltz, R.\,A., and Monteiro, P.\,B.,
\emph{High-order space--time finite element methods for the
Poisson--Nernst--Planck equations: Positivity and unconditional energy
stability},
Comput.\ Methods Appl.\ Mech.\ Engrg.\ \textbf{396}, 115070 (2022);
arXiv:2105.01163.
Extends the Metti et al.\ entropy/log-density formulation to arbitrary-order
space--time FEM with positivity preservation and unconditional energy
stability.

\end{thebibliography}

\end{document}
